\section{Actual conjugate compatibility in one quartic field}
\label{sec:single-quartic-compatible}

The splitting-field bound can be sharpened by retaining the fact that
the four linear factors come from one actual algebraic integer. Put
\[
 f(X)=X^4+3X^3+5X^2+3X+1,\qquad L=\mathbb Q(\alpha),\quad f(\alpha)=0,
\]
and let $K$ be its degree-eight splitting field. In this section $a,b$
are positive coprime integers and
\[
 F(a,b)=VQ^g,\qquad g\ge3,\qquad V,Q\ge1,\qquad V\text{ is }g\text{-free}.
\]
Shared support of $V$ and $Q$ is allowed. All heights are absolute
logarithmic Weil heights.

\begin{theorem}[A principal quartic root field]
The field $L$ has degree four, signature $(0,2)$, class number one, and
unit rank one. Its discriminant has absolute value at most $117$.
\end{theorem}
\begin{proof}
Over $\mathbb Q(\zeta)$ the two factors of $f$ are
$X^2+(2-\zeta)X+1$ and $X^2+(1+\zeta)X+1$. Choose $\alpha$ in the
first. Then $\zeta=\alpha+\alpha^{-1}+2$ and
$L=\mathbb Q(\zeta,\sqrt{-1-3\zeta})$. The norm of $-1-3\zeta$ is
$13$, so it is not a square in $\mathbb Q(\zeta)$. This proves the
degree and signature. The exact resultant calculation gives order
discriminant $117$, so the field discriminant is no larger and the
order index is supported only at three. Minkowski's ideal-class bound
\cite[Theorem~4.3]{MilneANT2020} gives every class an integral representative
of norm at most
\[
 \left(\frac4\pi\right)^2\frac{4!}{4^4}\sqrt{|\operatorname{Disc}L|}
 \le\frac{3\sqrt{117}}{2\pi^2}<\frac{11}{6}<2.
\]
The integral norm must be one, so every class is trivial. Dirichlet's
unit theorem gives rank one. Fix
$\mathcal O_L^*=\mu_L\langle\epsilon\rangle$ and $w_L=|\mu_L|$.
These classical number-field theorems are external inputs.
\end{proof}

\begin{theorem}[Extraction in the actual root field]
There are integral ideals $\mathfrak A,\mathfrak B$ with
\[
 (a-\alpha b)=\mathfrak A\mathfrak B^g,
 \qquad N\mathfrak A=V,\quad N\mathfrak B=Q.
\]
For fixed $V$, at most $4^{\omega(V)}$ choices for $\mathfrak A$ occur.
Choosing one integral generator $\delta$ of each such ideal, every actual
seed has an expression
\[
 a-\alpha b=\delta u\beta^g,
 \quad\beta\in\mathcal O_L,\quad
 u\in\{\xi\epsilon^k:\xi\in\mu_L,\ 0\le k<g\}.
\]
\end{theorem}
\begin{proof}
The norm of $a-\alpha b$ is $F(a,b)$. If $p\nmid117$ divides this
norm, then $b$ is nonzero modulo $p$. Every prime ideal dividing the
element has $\alpha\equiv a/b$ in its residue field. Factorization away
from the order index and the distinct roots of $f$ modulo $p$ show
that exactly one prime ideal above $p$ divides the element. Its residue
degree is one and its valuation is $v_p(F)$; there are at most four
possible such ideals as the seed varies.

At three no factor occurs. At thirteen the primitive exact-depth
theorem gives $v_{13}(F)=1$ whenever thirteen occurs. The norm formula
$\sum_{\mathfrak p\mid13}f_{\mathfrak p}v_{\mathfrak p}(a-\alpha b)=1$
again forces one degree-one prime of depth one. This argument makes no
unramifiedness assumption at thirteen. Write
$v_p(F)=r_p+gq_p$ with $r_p=v_p(V)<g$ and $q_p=v_p(Q)$, and allocate
these exponents at the unique prime above $p$. This proves the ideal
identities and the count. Principality supplies integral generators;
their quotient is a unit, which can be reduced modulo unit $g$th powers
by changing $\beta$.
\end{proof}

\begin{theorem}[A smaller compatible cover budget]
For fixed $V,g$, every actual seed lifts to one of at most
\[
 4w_L4^{\omega(V)}g
\]
covers over the fixed field $K$. Each is geometrically connected of
degree $g^3$ over $\mathbb P^1$ and genus $1+g^2(g-2)$. Its three
coefficients $\xi_i$ can be chosen with
\[
 h(\xi_i)\le\tfrac12\log V+C_Lg.
\]
For $\lambda=\max\{1,\log V\}/g$ the logarithm of the cover count
divided by $g$ is at most
\[
 2\lambda+\frac{\log g+\log(4w_L)}g.
\]
Allowing all integral $V\le e^{Ag}$ costs at most $4w_Lg e^{3Ag}$ covers.
\end{theorem}
\begin{proof}
Put $\kappa=\delta u$ and apply the four embeddings $\sigma_i:L\to K$.
All four equations are conjugates of the same equation. Their ratios give
\[
 y_i^g=\frac{x-\alpha_i}{\xi_i(x-\alpha_4)},\qquad
 \xi_i=\frac{\sigma_i(\kappa)}{\sigma_4(\kappa)},\quad
 x=a/b,\quad 1\le i\le3.
\]
The ideal choice and the single unit class determine all three ratios;
there are no three independent splitting-field unit choices. The looser
factor four in the stated count is harmless. Divisors at the four distinct
roots give Kummer independence also for composite $g$. Each of the four
inertia groups has order $g$, so Riemann--Hurwitz gives the stated genus.

Balance the two complex logarithmic absolute values of $\delta$ in a
fixed interval for the unit lattice. Their mean is $(\log V)/4$, and
integrality gives $h(\delta)\le(\log V)/4+C$. Since
$h(u)\le g h(\epsilon)$, the height inequalities for a ratio yield
$h(\xi_i)\le2h(\delta u)\le(\log V)/2+C_Lg$.
Finally $4^{\omega(V)}\le V^2$. Summing the resulting bound over at most
$e^{Ag}$ residual integers proves the union estimate, including the cost
of choosing $V$.
\end{proof}

\begin{theorem}[A rational complete-intersection model]
Each compatible class above has a smooth geometrically connected
complete-intersection model of degrees $(g,g)$ in $\mathbb P^3_{\mathbb Q}$,
capturing all its actual seeds. Its defining forms have logarithmic
projective coefficient heights $O(\log V+g)$, with fixed-field constants.
\end{theorem}
\begin{proof}
Write $\beta=X_0+X_1\alpha+X_2\alpha^2+X_3\alpha^3$. Expand
$\kappa\beta^g$ in this rational basis and set the coefficients of
$\alpha^2,\alpha^3$ equal to zero. These are two homogeneous degree-$g$
forms over $\mathbb Q$. Every actual seed supplies a rational point;
the reverse implication needs separate integral and primitive conditions.

Over an algebraic closure use the invertible embedding coordinates
$Z_i=\sigma_i(\beta)$. The image
$(\sigma_i(\kappa)Z_i^g)_i$ must lie on the projective line
$(a-\alpha_i b)_i$. The line meets the four coordinate hyperplanes at
distinct points, so at most one $Z_i$ vanishes. The intersection is
nonempty by taking $g$th roots of any line point. A dependence of the
two gradients would give a nonzero vector in the annihilator of the
line supported solely on a vanishing coordinate. No coordinate vanishes
identically on the line, so this is impossible. The intersection is
therefore smooth of pure dimension one. Kummer independence makes the
chart $Z_4\ne0$ connected; its complement is finite and contains no
hidden curve component. The genus is $1+g^2(g-2)$.

In the fixed basis, the rational coordinate height of $\kappa$ is
$O(\log V+g)$. Clear its common denominator first. The monic relation
for $\alpha$ bounds all multiplication coefficients in its powers
exponentially in $g$, while the multinomial coefficients have total
sum $4^g$. Thus both defining forms have the asserted coefficient height.
\end{proof}

These statements keep actual conjugate compatibility and improve the
preceding cover budget. They do not identify generic hard unit classes
with actual seeds. A small number of covers and controlled coefficient
heights still supply no uniform height bound for their individual points.
The actual lower bound
$h(a/b)\ge(g\log Q+\log V)/4-\log2$ is consistent with all these results.
The uniform actual-point and simultaneous-compression questions remain open.
